\FloatBarrier
\section{Gate Verification and Fidelity}
\label{app:gate-verification}

\label{sec:verification-benchmarking}

Section~\ref{sec:platform-readout-probabilities} explains why logical
populations alone do not determine a state or a gate. This appendix
develops the additional measurements that reconstruct states and
operations, and relates them to tests of fidelity and entanglement.

A population oscillation establishes one aspect of a gate, but it
does not fix its action on every input state. A phase error can be
invisible on $\ket0$ and $\ket1$ while changing their superpositions.
Three tests provide different information. State and operation
reconstruction determine what a gate does to selected inputs;
a Bell-state test checks whether a two-qubit gate has created
entanglement; random gate sequences test how errors accumulate.
They are alternatives for different questions, not three consecutive
steps required to execute a gate.

\subsection{Measuring a state beyond its 0 and 1 populations}
\label{sec:calibration-state-tomography}

The states $(\ket0+\ket1)/\sqrt2$ and
$(\ket0+i\ket1)/\sqrt2$ both give equal 0 and 1 populations.
A relative phase error between them is therefore invisible to
logical-basis readout. A known rotation before readout converts
the phase information into a population difference.
For a qubit that remains in its computational space, write
\begin{equation}
 \rho=\frac12(I+r_xX+r_yY+r_zZ).
 \label{eq:calibration-bloch-state}
\end{equation}
Logical-basis readout gives
$r_z=P_0-P_1$. To obtain $r_x$ or $r_y$, rotate that axis to the
readout axis before measuring. If $R$ is the analysis rotation,
the measured observable on the unrotated state is $R^\dagger ZR$:
\[
 \begin{array}{c|c}
 R & R^\dagger ZR\\ \hline
 I & Z\\
 R_y(-\pi/2)&X\\
 R_x(\pi/2)&Y
 \end{array}
\]
For a single spin, Section~\ref{sec:calibration-spin-rotations}
makes these analysis settings explicit. A resonant rectangular pulse of
duration $\pi/(2\Omega_0)$ with phase $+\pi/2$ gives $R_y(-\pi/2)$;
the same duration with phase zero gives $R_x(\pi/2)$.
Apply the analysis pulse while the electron remains isolated,
then switch to the fixed readout voltages.
Repeat preparation and the operation being tested for each
analysis setting in separate trials. Correct the outcome counts
using Eq.~\eqref{eq:calibration-assignment-map}.
The differences $P_0-P_1$ for the three settings give
$r_z,r_x,r_y$, and hence $\rho$; this is state \emph{tomography}.

For encoded qubits, the rotations act on the logical basis and
are generated by the exchange/gradient controls of
Section~\ref{sec:measurement-exchange-oscillations}. Once logical $x$ and $z$
rotations are available, the required $y$ rotation follows from
\[
 R_y(\vartheta)=R_z(\pi/2)R_x(\vartheta)R_z(-\pi/2),
\]
with the rightmost pulse applied first. This connects the analysis
axes to the same calibrated controls used for the gates.

This procedure reconstructs the state produced in one specified
experiment: three analysis settings give the three components of
its Bloch vector. It is complete as a state measurement.

\subsection{Reconstructing the action of a one-qubit gate}
\label{sec:calibration-process-tomography}

Different operations can give the same output from one input.
For example, an ideal $z$ rotation leaves the populations of
$\ket0$ unchanged whatever its angle, but changes the relative
phase of a superposition. To determine the gate rather than just
one output, repeat the state measurement for independent inputs. A trace-preserving
linear quantum operation acts on Bloch vectors as
\begin{equation}
 \mathbf r_{\mathrm{out}}=M\mathbf r_{\mathrm{in}}+\mathbf t,
 \label{eq:calibration-affine-channel}
\end{equation}
where $M$ is a real $3\times3$ matrix and $\mathbf t$ a real vector.
Start each trial from initialized $\ket0$. The preparations
\[
 \begin{array}{c|cccc}
 \text{input state}&\ket0&\ket1&\ket+&\ket{+i}\\ \hline
 \text{preparation rotation}&I&R_x(\pi)&R_y(\pi/2)&R_x(-\pi/2)
 \end{array}
\]
give four independent inputs, where
$\ket+=(\ket0+\ket1)/\sqrt2$ and
$\ket{+i}=(\ket0+i\ket1)/\sqrt2$.
Apply the same gate to each input and measure its three output
components by the analysis rotations above. Denote the output
vectors by $\mathbf r_0,\mathbf r_1,\mathbf r_+,\mathbf r_{+i}$.
Their input vectors are respectively
$\hat{\mathbf z},-\hat{\mathbf z},\hat{\mathbf x},\hat{\mathbf y}$,
so
\begin{equation}
 \begin{aligned}
 \mathbf t&=(\mathbf r_0+\mathbf r_1)/2,&
 M\hat{\mathbf z}&=(\mathbf r_0-\mathbf r_1)/2,\\
 M\hat{\mathbf x}&=\mathbf r_+-\mathbf t,&
 M\hat{\mathbf y}&=\mathbf r_{+i}-\mathbf t.
 \end{aligned}
 \label{eq:calibration-process-reconstruction}
\end{equation}
This determines the affine map from four input states, each measured
along three axes. It assumes that preparation and analysis rotations
are known. Without that assumption, the data characterize their
combined action with the gate rather than isolating the gate error.

\subsubsection{Interpreting a reconstructed operation}
The measurements have supplied $M$ and $\mathbf t$. We now compare
their entries with a specific error model to distinguish a coherent
phase error from loss of coherence and energy relaxation.
For example, remove the intended gate and suppose the remaining
operation consists of an unwanted $z$ rotation by $\chi$, random
$z$ phases, and relaxation $\ket1\to\ket0$ with probability
$\gamma$. The resulting map has the form
\begin{equation}
 M=\begin{pmatrix}
 \lambda\cos\chi&-\lambda\sin\chi&0\\
 \lambda\sin\chi& \lambda\cos\chi&0\\
 0&0&1-\gamma
 \end{pmatrix},\qquad
 \mathbf t=\begin{pmatrix}0\\0\\\gamma\end{pmatrix}.
 \label{eq:calibration-channel-diagnosis}
\end{equation}
To derive the last row, use
$\rho_{11}'=(1-\gamma)\rho_{11}$ and
$\rho_{00}'=\rho_{00}+\gamma\rho_{11}$.
For the first two rows, relaxation reduces $\rho_{01}$ by
$\sqrt{1-\gamma}$: its two Kraus operators are
$K_0=\ket0\bra0+\sqrt{1-\gamma}\ket1\bra1$ and
$K_1=\sqrt\gamma\ket0\bra1$, and substitution into
$\rho'=\sum_a K_a\rho K_a^\dagger$ gives this factor.
A phase rotation multiplies the off-diagonal entry by
$e^{-i\chi}$. An independent Gaussian phase with variance
$\sigma_\phi^2$ supplies the additional factor
$e^{-\sigma_\phi^2/2}$, so
$\lambda=\sqrt{1-\gamma}e^{-\sigma_\phi^2/2}$ in this model.

The four-state measurements now have a direct interpretation:
$\gamma=t_z$ describes motion of the Bloch sphere toward $\ket0$,
$\lambda=\sqrt{M_{xx}^2+M_{yx}^2}$ describes transverse contraction,
and $\chi=\operatorname{atan2}(M_{yx},M_{xx})$ gives the residual
phase. A phase correction can remove $\chi$; it cannot reverse
the contraction or restore the lost excitation. An observed map
that lacks this structure requires a different error model,
rather than forcing all deviations into one rotation angle.

\subsection{Comparing the measured operation with the target gate}
\label{sec:calibration-gate-fidelity}

The reconstructed map predicts an output for every input qubit state.
To quantify its discrepancy from the requested gate, compare those
outputs with the ideal ones before averaging over inputs. This
distinguishes an error measure from a parameter such as a Rabi
contrast or a readout misclassification rate.

\subsubsection{Averaging the output overlap}
Let the target unitary $U$ rotate Bloch vectors by the real matrix
$R_U$. For a pure input with unit vector $\mathbf r$, the overlap
of the measured output with the ideal output is
\[
 F(\mathbf r)=
 \frac{1+(R_U\mathbf r)\cdot(M\mathbf r+\mathbf t)}2.
\]
Averaging uniformly over input directions uses
$\langle r_i\rangle=0$ and
$\langle r_i r_j\rangle=\delta_{ij}/3$, giving
\begin{equation}
 F_{\mathrm{avg}}=\frac12+
            \frac{\operatorname{Tr}(R_U^T M)}6.
 \label{eq:calibration-average-fidelity}
\end{equation}
This is an average over specified input states, distinct from the
noise average that may already be included in the operation.
For a coherent angle error $\eta$ about the correct axis,
$\operatorname{Tr}(R_U^TM)=1+2\cos\eta$, so
$1-F_{\mathrm{avg}}=(1-\cos\eta)/3\simeq\eta^2/6$.

\subsubsection{Accounting for population that leaves the qubit space}
The preceding comparison assumed a trace-preserving map within a
two-dimensional space. If a pulse also excites an orbital, valley,
or excluded spin state, counting only the surviving qubit outcomes
omits part of the error. For the computational-space projector
$P_{\mathrm q}$, determine or bound
\begin{equation}
 L(\rho)=1-\operatorname{Tr}[P_{\mathrm q}\mathcal E(\rho)],
 \label{eq:calibration-leakage-probability}
\end{equation}
where $\mathcal E$ is the physical operation on the full state space.
Renormalizing the surviving qubit population hides $L$.
A measurement that labels every triplet as logical 1, for example,
cannot distinguish a logical triplet from a leaked spin state.
An additional state-selective measurement is needed before reporting
leakage separately. The trace-preserving qubit formula above applies
when such leakage is negligible.

\subsection{Testing whether a two-qubit gate creates entanglement}

Use the calibrated conditional-phase interaction to prepare an
entangled state. After its local phases are corrected, a CZ gives
CNOT by a basis change on qubit $B$:
\[
 \mathrm{CNOT}_{A\to B}
 =(I_A\otimes\mathsf H_B)\,\mathrm{CZ}\,
                   (I_A\otimes\mathsf H_B),
 \qquad \mathsf H=(X+Z)/\sqrt2.
\]
The Hadamard $\mathsf H$ can be implemented, up to global phase,
as $R_y(\pi/2)R_z(\pi)$. Initialize both qubits in 0, prepare
$\ket+_A$ with $R_{y,A}(\pi/2)$, and apply this CNOT sequence.
Ideally the output is
$\ket{\Phi^+}=(\ket{00}+\ket{11})/\sqrt2$.
However, observing equal probabilities of $00$ and $11$ does not
establish this coherent state: the mixture
$(\ket{00}\bra{00}+\ket{11}\bra{11})/2$ gives exactly those
populations. We therefore measure correlations in rotated bases as well as these
populations, using the one-qubit analysis pulses already established.

\subsubsection{The three measurements required by this test}
The Bell-state projector has the Pauli expansion
\[
 \ket{\Phi^+}\bra{\Phi^+}
 =\frac{I+X_AX_B-Y_AY_B+Z_AZ_B}{4}.
\]
For the measured output $\rho$, its overlap with the target is
therefore
\begin{equation}
 F_{\Phi^+}
 =\frac{1+\langle X_AX_B\rangle-\langle Y_AY_B\rangle
                    +\langle Z_AZ_B\rangle}{4}.
 \label{eq:calibration-bell-test}
\end{equation}
Measure $Z_AZ_B$ by joint logical-basis readout:
$\langle Z_AZ_B\rangle=P_{00}+P_{11}-P_{01}-P_{10}$.
To measure $X_AX_B$, apply $R_y(-\pi/2)$ to both qubits before
the same readout. Applying $R_x(\pi/2)$ to both measures
$Y_AY_B$. The three measurement settings thus determine the
displayed fidelity without reconstructing all entries of $\rho$.
If the two assignments are independent, the joint assignment
matrix is the tensor product of the individual matrices in
Eq.~\eqref{eq:calibration-assignment-map}; otherwise joint
reference preparations are needed.

\subsubsection{Why the result can establish entanglement}
For any normalized product state $\ket a\ket b$,
\[
 |\langle\Phi^+|a,b\rangle|^2
 =\frac{|a_0b_0+a_1b_1|^2}{2}\le\frac12,
\]
by the Cauchy--Schwarz inequality. Convex mixtures of product
states obey the same bound. A measured
$F_{\Phi^+}>1/2$, with the measurement uncertainties accounted
for, therefore establishes entanglement in this output.
The incoherent $00/11$ mixture has correlations $(0,0,1)$
and gives exactly $F_{\Phi^+}=1/2$.
This tests one important action of the entangling gate; it does
not determine the gate on every input state.

\subsection{Accumulating error with random gate sequences}

To detect small gate errors, apply many gates whose ideal product
is known to return the state to its starting point. Individual
sequences can amplify or cancel coherent errors differently.
Averaging over random sequences reduces dependence on that
particular choice.

\subsubsection{Constructing and measuring a sequence with known ideal output}
For a one-qubit protocol, choose gates from the \emph{Clifford group}:
these unitaries permute the Pauli axes, with possible sign changes.
For each sequence, draw $m$ independent uniformly random Cliffords
$C_1,\ldots,C_m$, and append their ideal inverse
\[
 C_{\mathrm{inv}}=(C_m\cdots C_1)^\dagger.
\]
Compile each Clifford into the calibrated logical $\pi/2$ and
$\pi$ rotations. For example,
$C=R_y(\pi/2)R_x(\pi/2)$ is made by an $x$ pulse followed by a
$y$ pulse. For a single spin, their phases are zero and $-\pi/2$;
both have duration $\pi/(2\Omega_0)$. Compute $C_{\mathrm{inv}}$
as one Clifford and compile it using the same set of short
implementations, rather than physically reversing all $m$ gates.
Thus increasing $m$ increases the tested sequence length without
also giving the final inverse a length proportional to $m$.

Prepare $\ket0$, run the sequence, and use the fixed readout
procedure to count outcomes 0. Average this return probability
over independently drawn sequences at each length $m$.
In the ideal experiment every sequence returns to 0; errors in
the implemented pulses reduce the return probability. This is
\emph{randomized benchmarking}
\cite{Magesan2011}.

\subsubsection{Deriving the dependence on sequence length}
The measurement supplies a return probability at each length.
To infer an error per Clifford from that curve, we need a noise
model. The simplest decay model assumes the same trace-preserving noise
channel $\mathcal N$ after each Clifford, negligible leakage,
and noise without memory between gates. We can derive the decay
using the Bloch representation
$\mathbf r\mapsto M_{\mathcal N}\mathbf r+\mathbf t_{\mathcal N}$.
Changing the axes by a Clifford rotation $R_C$ gives the noise
matrix $R_C^T M_{\mathcal N}R_C$ and translation
$R_C^T\mathbf t_{\mathcal N}$.

First average over the four Pauli conjugations. Their Bloch
rotations are $I$, $\operatorname{diag}(1,-1,-1)$,
$\operatorname{diag}(-1,1,-1)$ and
$\operatorname{diag}(-1,-1,1)$. Every translation component
and every off-diagonal matrix entry occurs with opposite signs
and cancels. Averaging the remaining diagonal matrix over
Clifford permutations of $x,y,z$ replaces all three diagonal
entries by their mean. Hence
\begin{equation}
 \frac1{24}\sum_C R_C^T M_{\mathcal N}R_C=pI,\qquad
 p=\frac{\operatorname{Tr}M_{\mathcal N}}3,\qquad
 \frac1{24}\sum_C R_C^T\mathbf t_{\mathcal N}=0.
 \label{eq:calibration-clifford-average}
\end{equation}
There are 24 distinct single-qubit Clifford rotations after
discarding global phases.

Why does this average appear in the pulse sequence?
Write $D_j=C_j\cdots C_1$. Inserting $D_jD_j^\dagger$
between the errors and using the final ideal inverse puts the
sequence's error into the form
\[
 \mathcal N\circ
 ( \mathcal D_m^{-1}\circ\mathcal N\circ\mathcal D_m)
 \circ\cdots\circ
 ( \mathcal D_1^{-1}\circ\mathcal N\circ\mathcal D_1),
\]
where $\mathcal D_j(\rho)=D_j\rho D_j^\dagger$, $D_0=I$, and the leftmost
$\mathcal N$ is the final-inverse error.
The uniformly and independently selected $C_j$ make the
cumulative $D_j$ uniformly and independently distributed as
well: the map between the two lists is invertible,
$C_j=D_jD_{j-1}^\dagger$. Each parenthesis therefore averages
to the same isotropic channel~\cite{Magesan2012}.
In density-matrix notation, that channel and its repeated action are
\begin{equation}
 \mathcal D_p(\rho)=p\rho+(1-p)\frac I2,\qquad
 \mathcal D_p^m(\rho)=p^m\rho+(1-p^m)\frac I2.
 \label{eq:calibration-rb-channel}
\end{equation}
With perfect preparation and readout, the return probability is
$(1+p^m)/2$. Fixed preparation, readout and final-inverse errors
change the constants, giving the fit
\begin{equation}
 \overline P_{\mathrm{return}}(m)=A p^m+B.
 \label{eq:calibration-rb-fit}
\end{equation}
\subsubsection{Extracting and interpreting the error}
Fit $A,B,p$ from the length scan. For a pure input,
$\bra\psi\mathcal D_p(\ket\psi\bra\psi)\ket\psi=(1+p)/2$;
the corresponding mean infidelity per Clifford is therefore
\begin{equation}
 r_{\mathrm C}=\frac{1-p}{2}.
 \label{eq:calibration-rb-infidelity}
\end{equation}
For three equally spaced lengths $m,m+d,m+2d$, the exact
exponential model implies
\[
 \frac{\overline P(m+2d)-\overline P(m+d)}
      {\overline P(m+d)-\overline P(m)}=p^d.
\]
The unknown offset $B$ and amplitude $A$ cancel. This shows how
a length scan distinguishes gate accumulation from fixed readout
contrast. In data analysis, use many lengths and random sequences
in a common fit; the three-point ratio becomes unstable when the
differences are comparable to counting noise.

For example, a fixed small rotation error $R_z(\eta)$ after
each ideal Clifford has
$p=(1+2\cos\eta)/3\simeq1-\eta^2/3$.
Randomization gives
$p^m\simeq e^{-m\eta^2/3}$.
Repeating the same unrandomized rotation error instead produces
$R_z(m\eta)$ and an equatorial-state error
$\sin^2(m\eta/2)\simeq m^2\eta^2/4$.
Thus a good mean benchmark does not make a coherent error harmless
in every circuit. The signed calibration signals identify such
an error; benchmarking summarizes its accumulation under the
specified randomization.

The interpretation of $r_{\mathrm C}$ relies on the noise model. Strong drift,
gate dependence or leakage can invalidate this simple exponential
description~\cite{Magesan2012}. Also, a Clifford may require several
physical pulses, so $r_{\mathrm C}$ is not automatically the error
of each pulse. Tomography identifies how an operation acts on states;
benchmarking tests how errors accumulate in a specified gate set.
